%~Mouliné par MaN_auto v.0.37.1 (a7a71208) 2025-10-24 11:26:47
\documentclass[AHL,Unicode]{cedram}


\usepackage{aligned-overset}
\DeclareMathOperator{\Exp}{Exp}
\DeclareMathOperator{\supp}{supp}
%\DeclareMathOperator{\dt}{dt}

\renewcommand*\r{\mathbb{R}}
\newcommand*\n{\mathbb{N}}
\newcommand*\z{\mathbb{Z}}
\newcommand*\one{\mathbf{1}}
\newcommand*\p{\mathbb{P}}
\newcommand*\all{\forall}
\newcommand*\eps{\varepsilon}
\newcommand*\tensor{\otimes}
\newcommand*\indexx[2]{n\left(#1,#2\right)}
\newcommand*\nonnegative{\left[0,\infty\right)}
\newcommand*\cal[1]{\mathcal{#1}}
\newcommand*\union{\cup}
\newcommand*\abs[1]{\left|#1\right|}

\newcommand*\intersect{\cap}


\def\hypergeometric#1#2#3#4{_{2}F_{1}\left(#1,#2;#3;#4\right)}
%%%%%%%%%%%%%%%%%%%%%%%%%%%%%%%%%%%%%%%%%%%%%%%%%%%%%%%%
%~ This fixes spacing issues around \left( and \right) and other similar symbols
\let\originalleft\left 
\let\originalright\right
\renewcommand{\left}{\mathopen{}\mathclose\bgroup\originalleft}
\renewcommand{\right}{\aftergroup\egroup\originalright}
%%%%
\renewcommand*{\to}{\mathchoice{\longrightarrow}{\rightarrow}{\rightarrow}{\rightarrow}}
%%%%
\let\oldin\in
\renewcommand{\in}{\mathchoice{\oldin}{\oldin}{\,\oldin\,}{\,\oldin\,}}

%%%%%%%%%%%%%%%%%%%%%%%%%%%%%%%%%%%%%%%%%%%%%%%%%%%%%%%%%%%%
\graphicspath{{./figures/}}

\newcommand*{\mk}{\mkern -1mu}
\newcommand*{\Mk}{\mkern -2mu}
\newcommand*{\mK}{\mkern 1mu}
\newcommand*{\MK}{\mkern 2mu}

\hypersetup{urlcolor=purple, linkcolor=blue, citecolor=red}

\newcommand*{\relabel}{\renewcommand{\labelenumi}{(\theenumi)}}
\newcommand*{\romanenumi}{\renewcommand*{\theenumi}{\roman{enumi}}\relabel}
\newcommand*{\Romanenumi}{\renewcommand*{\theenumi}{\Roman{enumi}}\relabel}
\newcommand*{\alphenumi}{\renewcommand*{\theenumi}{\alph{enumi}}\relabel}
\newcommand*{\Alphenumi}{\renewcommand*{\theenumi}{\Alph{enumi}}\relabel}
\let\oldtilde\tilde
\renewcommand*{\tilde}[1]{\mathchoice{\widetilde{#1}}{\widetilde{#1}}{\oldtilde{#1}}{\oldtilde{#1}}}
\let\oldforall\forall
\renewcommand*{\forall}{\mathrel{\oldforall}}
%%%%%%%%%%%%%%%%%%%%%%%%%%%%%%%%%%%%%%%%%%%%%%%%%%%%%%%%%%%%%%%%%%%%%

\title[The distance problem via subadditivity]{The distance problem via subadditivity}
\alttitle{Le problème de la distance via la sous-additivité}

\subjclass{28C99, 51F99, 51K99, 60D05}
\keywords{Measured metric space, trees, subadditivity, metric preserving function}



\author[\initial{R.} \lastname{Gross}]{\firstname{Renan} \lastname{Gross}}
\address{University of Cambridge,\\
Faculty of Mathematics,\\
Wilberforce Road,\\
Cambridge CB3 0WA (United Kingdom)}
\email{rg751@cam.ac.uk}


\begin{abstract}
In a recent paper, Aldous, Blanc and Curien asked which distributions can be expressed as the distance between two independent random variables on some separable measured metric space. We show that every nonnegative discrete distribution whose support contains $0$ arises in this way, as well as a class of compactly supported distributions with density.
\end{abstract}

\begin{altabstract}
Dans un article récent, Aldous, Blanc et Curien ont demandé quelles distributions peuvent s'écrire comme la distance entre deux variables aléatoires indépendantes sur un espace métrique mesuré séparable. Nous montrons que toute distribution discrète positive ou nulle dont le support contient $0$ apparaît de cette manière, ainsi qu'une classe de distributions à densité à support compact.
\end{altabstract}





\datereceived{2025-01-22}
\daterevised{2025-07-11}
\dateaccepted{2025-07-18}

\editors{S.~Gou\"ezel and D.~Aldous}
%%%%%%%%%%%%%%%%%%%%%%%%%%%%%%%%%%%%%%%%%%%%%%%%%%%%%%%%%%%%%%%%%
\begin{document}
\maketitle

\section{Introduction}
\subsection{The problem}

The \emph{distance problem} asks the following. Given a distribution $\theta$ on $\nonnegative$, is it possible to find a complete separable metric space $(S,d)$ and a Borel probability measure $\mu$ on $S$, so that when $X,Y\sim\mu$ are independent, their distance satisfies $d(X,Y)\sim\theta$? If so, we say that $(S,d,\mu)$
\emph{achieves} $\theta$, and that $\theta$ is \emph{feasible}. This very natural and quite inviting problem was first proposed by Aldous, Blanc and Curien in~\cite{aldous_blanc_curien_distance_problem_2024}, where they also give some statistical motivation.

\subsection{Examples}



While any $(S,d,\mu)$ gives some distribution $\theta$ which can in principle be calculated, the inverse problem, of finding an explicit space which achieves a given distribution, is harder. The following examples show how to achieve some common distributions, and may give some intuition about the types of constructions one can use.
\begin{enumerate}%[left=0pt]
\item \emph{The uniform measure:} let $S=S^{1}=\r\mod\z$ be the unit circle, $\mu$ the uniform measure on $S^{1}$, and $d(x,y)$ the arc length between $x$ and $y$. Then $d(X,Y)\sim U[0,\frac{1}{2}]$.
\item \emph{The exponential distribution:} let $S=\r$, let $d(x,y)=\abs{x-y}$, and let $\mu=
(1)$ be the exponential distribution. By the memoryless property, $\abs{X-Y}$ is also exponentially distributed.
\item \emph{Bernoulli distribution:}\label{enu:bernoulli_construction} let $0< p_{0} < 1$. In order to achieve the Bernoulli distribution $\theta=p_{0}\delta_{0}+(1-p_{0})\delta_{1}$, let $m\geq2$ be an integer and let $\alpha\in[0,1]$ to be chosen later. Let $S=\{ 0,\ldots,m-1\} $, equipped with the discrete metric
\[
d\left(x,y\right)=
\begin{cases}
0 & x=y\\
1 & x\neq y.
\end{cases}
\]
Let $\mu$ be the distribution of the random variable on $S$ whose value is $0$ with probability $\alpha$, and uniform among $\{ 1,\ldots,m-1\} $ with probability $1-\alpha$. Then
\begin{equation}\label{eq:bernoulli_probability}
\p\left[d\left(X,Y\right)=0\right]=\alpha^{2}+\left(\frac{1-\alpha}{m-1}\right)^{2}.
\end{equation}
Choosing $m$ so that the quadratic equation $p_{0}=\alpha^{2}+(\frac{1-\alpha}{m-1})^{2}$ has a solution and solving for $\alpha$ gives us the desired probability space.
\item \emph{The uniform measure, revisited:} let $S=\{ 0,1\} ^{\n}$, let $\mu$ be the measure on $S$ where all entries are i.i.d. $0$-1 Bernoulli random variables with success probability $1/2$, and let $d(x,y)=\sum_{k=1}^{\infty}2^{-k}\abs{x_{i}-y_{i}}$. Then $d(X,Y)\sim U[0,1]$. This is essentially a more explicit version of the first example.
\item \emph{The exponential distribution, revisited:} let $S=(0,1]$, $\mu$ the uniform measure on $S$, and $d(x,y)=\lvert\int_{x}^{y}\frac{1}{t}d t\rvert=\lvert\log\frac{y}{x}\rvert$. A short calculation shows that $\p[d(X,Y)\leq t]=1-e^{-t}$, and so $d(X,Y)\sim\Exp(1)$.
\item \emph{Non-example:} let $\theta$ be any distribution with density on $\nonnegative$, and denote its cumulative distribution function (CDF) by $F$. Let $S=\nonnegative$, let
\[
d\left(x,y\right)=
\begin{cases}
0 & x=y\\
\max\left(x,y\right) & x\neq y,
\end{cases}
\]
and let $\mu$ be the distribution whose CDF is $\sqrt{F}$. Since $\theta$ has density, the event $\{ x=y\} $ has measure $0$ under $\mu$, and so $d(X,Y)=\max(X,Y)$ with probability~$1$. The CDF of the maximum of two i.i.d. random variables is the square of their individual CDF, and so $d(X,Y)\sim\theta$. However, while $\mu$ is a Borel measure with respect to the Euclidean topology, it is not a Borel measure with respect to the topology induced by $d$. Note also that $(S,d)$ is not separable (for small enough~$\eps$, the only element $\eps$-close to $x$ is $x$ itself).
\end{enumerate}

\subsection{Results}

In~\cite{aldous_blanc_curien_distance_problem_2024}, Aldous, Blanc and Curien showed some basic constraints on the support of the distribution (for example, it must contain $0$) and gave several constructions and approximations. In their main constructions, the metric space $S$ was a weighted tree and the measure $\mu$ a distribution on its leaves/rays. Choosing the weights and distribution in a clever way, they showed that every finite discrete distribution with $0$ in its support is feasible. However, when applied to non-finite distributions via finite approximations, their approach yields non-separable spaces.

Our first result is a new construction which does not involve taking limits of finite distributions, avoiding limiting questions of separability. This result resolves the question of discrete distributions.

\begin{theo}\label{thm:discrete_distributions}
Every nonnegative discrete distribution $\theta$ supported on $0$ is feasible. \hyperlink{target:proof_of_discrete}{({proof $\nearrow$})}%revoir la destination ? 
\end{theo}

Rather than building a metric space $(S,d,\mu)$ which directly achieves $\theta$ from scratch, our proof combines together several simpler spaces $(S_{n},d_{n},\mu_{n})$, allowing us to sample from different $\theta_{n}$s with different probabilities. The main challenge is to decompose $\theta$ into subdistributions so that combining the spaces preserves the triangle inequality. As will be described in the next section, this construction can also be interpreted as a metric on the rays of a tree, where two rays are close if they travel together for a long time, generalizing a commonly used metric on the space of ends of a tree.

Our second result concerns distributions with density.
\begin{theo}\label{thm:bounded_continuous_distributions}
Let $g$ be a probability density function on $[0,1]$, and suppose that there exist constants $c,C>0$ such that
\begin{equation}\label{eq:bounds_on_density_function}
c<g(t)<C
\end{equation}
for all $t\in[0,1]$. Then the distribution $\theta$ with density $g$ is feasible. \hyperlink{target:proof_of_continuous}{({proof $\nearrow$})}
\end{theo}

The theorem is derived by altering the metric of a known-to-be-feasible distribution; here, subadditivity is used in order to preserve the metric. The theorem represents one step towards answering the open problem in~\cite{aldous_blanc_curien_distance_problem_2024}.

\begin{enonce}
{Problem}
Prove that for every probability density function $f$ on $\nonnegative$ whose support contains $0$, the distribution with density $f$ is feasible.
\end{enonce}
\subsection*{Acknowledgments}
We thank Wendelin Werner and Guillaume Blanc for useful comments, and the anonymous referee for their corrections and suggestions.
\pagebreak
\section{Proof of Theorem~\ref{thm:discrete_distributions}}

\hypertarget{target:proof_of_discrete}Our proof uses the following construction, which, given several feasible distributions, returns a mixture of these distributions. The resultant distance can be thought of as being computed by a ``selection'' mechanism: we linearly order the distributions, and iteratively flip a coin for each one to see whether or not it is the one that we sample from; if none of the distributions are chosen, the default $0$ is returned.

\subsection{Selection space}

Denote the integers by $\z$ and the positive integers by $\n=\{ 1,2,\ldots\} $. We say that a triplet $(S,d,\mu)$ is a measured metric space if $(S,d)$ is a metric space and $\mu$ is a Borel probability measure on $S$ with respect to $d$.

Let $I\subseteq\z$, and for each $n\in I$, let $(S_{n},d_{n},\mu_{n})$ be a complete countable measured metric space with more than $1$ point. Assume that for every $n\in I\intersect\n$ there exists a special $s_{n}\in S_{n}$ such that
\begin{equation}\label{eq:prob_sums_are_finite}
\sum_{n\in I\,\intersect\,\n}1-\mu_{n}\left(\left\{ s_{n}\right\} \right)<\infty.
\end{equation}
We now define a new space $(S,d,\mu)$ out of these spaces; later, in Proposition~\ref{prop:selection_space_properties}, we give conditions under which it is a complete separable measured metric space.
\begin{enumerate}
\item \emph{The space}: we say that a vector $x\in\prod_{n\in I}S_{n}$ is \emph{eventually constant} if there exists an $N>0$ such that $x_{n}=s_{n}$ for all $n\geq N$ (we still index the vector $x$ by the set $I$). Note that when $\sup I<\infty$, every vector is eventually constant. We set
\[
S=\left\{ x\in\prod_{n\in I}S_{n}\,\middle| \,x\text{ is eventually constant}\right\}.
\]
\item \emph{The function }$d$: for $x\neq y\in S$, let $\indexx xy=\sup\{ n\mid x_{n}\neq y_{n}\} $. This is always finite since $x$ and $y$ are eventually constant. We then define
\[
d\left(x,y\right)=
\begin{cases}
0 & x=y\\
d_{\indexx xy}\left(x_{\indexx xy},y_{\indexx xy}\right) & x\neq y.
\end{cases}
\]
\item \emph{The measure}: let $\tilde{\mu}=\tensor_{n\in I}\mu_{n}$ be the product measure on $\prod_{n\in I}S_{n}$. By~\eqref{eq:prob_sums_are_finite} the sum $\sum_{n\in I\intersect\n}1-\mu_{n}(\{ s_{n}\} )$ is finite, and the Borel--Cantelli lemma implies that $X\sim\tilde{\mu}$ is eventually constant with probability $1$ under $\tilde{\mu}$. The measure $\tilde{\mu}$ can therefore be restricted to a measure $\mu$ on $S$.
\end{enumerate}
We call the tuple $(S,d,\mu)$ the \emph{selection space} of $\{ (S_{n},d_{n},\mu_{n})\} _{n\in I}$. Let us find the distribution of $d(X,Y)$ when $X,Y\sim\mu$ are independent. The probability of obtaining $0$ is
\begin{equation}\label{eq:prob_of_0}
\p\left[d\left(X,Y\right)=0\right]=\prod_{n\in I}\p\left[X_{n}=Y_{n}\right]=\prod_{n\in I}\p\left[d_{n}\left(X_{n},Y_{n}\right)=0\right],
\end{equation}
while for $r>0$,
\begin{multline}\label{eq:prob_of_positive}
\p\left[d\left(X,Y\right)=r\right] \\
\begin{aligned}
& =\sum_{n\in I}\p\left[n\left(X,Y\right)=n\right]\cdot\p\left[d_{n}\left(X_{n},Y_{n}\right)=r\,\middle|\,\indexx XY=n\right]\\
& =\sum_{n\in I}\left(\prod_{k\in I,\,k\,>\,n}\p\left[X_{k}=Y_{k}\right]\right)\p\left[X_{n}\neq Y_{n}\right]\cdot\p\left[d_{n}\left(X_{n},Y_{n}\right)=r\,\middle|\, X_{n}\neq Y_{n}\right] \\
& =\sum_{n\in I}\left(\prod_{k\in I,\,k\,>\,n}\p\left[d\left(X_{k},Y_{k}\right)=0\right]\right)\p\left[d_{n}\left(X_{n},Y_{n}\right)=r\right],
\end{aligned}
\end{multline}
and we have indeed obtained a mixture of positive elements of the distributions achieved by $(S_{n},d_{n},\mu_{n})$.
\begin{rem}
As a space of vectors indexed by $I\subseteq\z$, the selection space can also be interpreted as the steps of a bi-directional walk, where at time-step $n$, the walker may go in the ``direction'' $x_{n}\in S_{n}$. When $\sup I$ is finite and we view time as starting at $\sup I$ and directed down towards $\inf I$, the index $\indexx xy$ is the first time that two walkers $x$ and $y$ split up when starting at a common origin. In this case the elements of $S$ can be seen as the space of ends on some rooted tree, and the function $d$ is a generalization of the metric $e^{-h(x,y)}$ on the space of ends, where $h(x,y)$ is the depth of the last common vertex of the rays $x$ and $y$. See Figure~\ref{fig:space_of_ends_metric}. When $\sup I=\infty$, however, there is no such common root.

\begin{figure}[h]
\begin{centering}
\includegraphics[scale=0.75]{space_of_ends_metric}
\end{centering}
\caption{\label{fig:space_of_ends_metric}
When $\sup I=N<\infty$, the function $d$ is a generalization of a class of metrics on the space of ends of a tree rooted, where the distance between two rays from the origin depends on the depth of their last common vertex.}
\end{figure}
\end{rem}

\begin{prop}\label{prop:selection_space_properties}
Let $I\subseteq\z$ and $\{ (S_{n},d_{n},\mu_{n})\} _{n\in I}$ be as above. Assume that the following two properties hold.
\begin{enumerate}
\item\label{prop2.2.1} For every $m<n\in I$, every $x_{1}\neq x_{2}\in S_{m}$ and every $y_{1}\neq y_{2}\in S_{n}$,
\begin{equation}\label{eq:metrics_are_subadditive_assumption}
d_{m}\left(x_{1},x_{2}\right)\leq2d_{n}\left(y_{1},y_{2}\right).
\end{equation}
\item\label{prop2.2.2} We have
\begin{equation}\label{eq:assumption_separability}
\text{either }
\inf I>-\infty\quad
\text{or}\quad
\lim_{n\to-\infty}\,\sup_{z,w\in S_{n}}d_{n}\left(z,w\right)=0.
\end{equation}
\end{enumerate}
Then the selection space $(S,d,\mu)$ is a complete separable measured metric space.
\end{prop}

\begin{proofc}
\begin{proof}[{Metric}]
It is clear that $d(x,y)=0$ if and only if $x=y$, so we only need to verify the triangle inequality. Let $x,y,z\in S$ be distinct, and denote $b=\indexx xy$.
\begin{enumerate}
\item If either $\indexx yz>b$ or $\indexx xz>b$, then since $x$ and $y$ agree on indices greater than $b$, we must have $\indexx xz=\indexx yz>b$, and also $d(x,z)=d(y,z)$. Then by~\eqref{eq:metrics_are_subadditive_assumption},
\[
d\left(x,y\right)=d_{b}\left(x_{b},y_{b}\right)\leq2d_{\indexx xz}\left(x_{\indexx xz},z_{\indexx xz}\right)=d\left(x,z\right)+d\left(y,z\right).
\]
\item If both $n(y,z)=b$ and $n(x,z)=b$, then the triangle inequality is satisfied because $d_{b}$ is a metric.
\item Otherwise, we have $n(y,z),n(x,z)\leq b$, and at least one of $n(y,z)$ or $n(x,z)$ is strictly smaller than $b.$ But it is impossible for both to be strictly smaller, since that would mean that $x_{b}=z_{b}=y_{b}$, contradicting the fact that $x_{b}\neq y_{b}$. Supposing w.l.o.g. that $\indexx xz=b$, we then trivially have
\[
d\left(x,y\right)=d_{b}\left(x_{b},y_{b}\right)=d_{b}\left(x_{b},z_{b}\right)=d\left(x,z\right)\leq d\left(x,z\right)+d\left(y,z\right).
\]
\end{enumerate}
\let\qed\relax
\end{proof}
\begin{proof}[{Completeness}]
Let $N\in I$, and let $a\neq b\in S_{N}$. Then by~\eqref{eq:metrics_are_subadditive_assumption}, for every $n>N$ and every $w,z\in S_{n}$, $d_{n}(w,z)\geq\frac{1}{2}d_{N}(a,b)$. Thus, for every $N\in I$,
\begin{equation}
\inf_{n>N}\inf_{w,z\in S_{n}}d_{n}\left(w,z\right)>0.\label{eq:inf_on_tail}
\end{equation}
Let $(x^{(k)})_{k\in\n}$ be a Cauchy sequence in $S$, and let $L=\inf I$. When $L=-\infty$, \eqref{eq:inf_on_tail} implies that if $d(x^{(k_{1})},x^{(k_{2})})<\eps$ for some $k_{1}$ and all $k_{2}\geq k_{1}$ and small enough $\eps$, then $x_{n}^{(k_{2})}=x_{n}^{(k_{1})}$ for all $n>N$ for some $N\in I$, with $N\to-\infty$ as $\eps\to0$. Thus $x^{(k)}$ converges to an element in $S$. When $L>-\infty$, \eqref{eq:inf_on_tail} implies that for small enough $\eps$, $x_{n}^{(k_{2})}=x_{n}^{(k_{1})}$ for all $n>L$, while the first elements $(x_{L}^{(k)})_{k\in\n}$ form a Cauchy sequence in $S_{L}$. Since $S_{L}$ is itself complete, we again get that $x^{(k)}$ converges to an element in~$S$.\let\qed\relax
\end{proof}
\begin{proof}[{Separability}] 
If $\inf I>-\infty$, then $S$ is a countable union of countable sets, and so is countable itself. Otherwise, by~\eqref{eq:assumption_separability}, $\lim_{n\to-\infty}\sup_{z,w\in S_{n}}d(z,w)=0.$ For each negative $n\in I$, let $a^{(n)}\in S_{n}$, and set
\[
A_{n}=\left\{ x\in S\,\middle|\, x_{i}=a^{\left(i\right)}\,\all i\leq n\right\}.
\]
Since all vectors in $S$ are eventually constant, each $A_{n}$ is a countable union of countable sets and is therefore countable, and so the union $\union_{n\in I}A$ is also countable. It is also dense in $S$: for any $x\in S$ and any $n\in I$, the set $A_{n}$ contains an element $y$ such that $y_{k}=x_{k}$ for all $k>n$, and so $n(x,yt)\leq n$; denseness follows since $\sup_{z,w\in S_{n}}d_{n}(z,w)\to0$ as $n\to-\infty$.\let\qed\relax
\end{proof}

\begin{proof}[{Borel measure}]
Under the metric $d$, an open ball around $x\in S$ of radius $r$ is given by
\[
B\left(x,r\right)=\bigcup_{N\in I}\quad\bigcup_{s\in S_{N}\mid d_{N}\left(x_{N},s\right)<r,s\neq x_{N}}\left\{ y\in S\,\middle|\, y_{N}=s,y_{n}=x_{n}\text{ for all }n>N\right\} \union\left\{ x\right\}.
\]
Each set of the form $\{ y\in S\mid y_{N}=s,y_{n}=x_{n}\text{ for all } n>N\} $ is measurable under the product measure $\tilde{\mu}$, and therefore so is the countable union of such sets. The restriction measure $\mu$ is therefore defined on the $\sigma$-algebra generated by the open balls in $S$. Since $(S,d)$ is separable, $\mu$ is Borel.
\end{proof}\let\qed\relax
\end{proofc}
\begin{prop}\label{prop:positive_interval_is_feasible}
Let $a>0$, let $p_{0}\in(0,1)$
and let $\lambda$ be a discrete probability distribution with $\supp(\lambda)\subseteq[a,2a].
$
Then the distribution $\theta=p_{0}\delta_{0}+(1-p_{0})\lambda$ is feasible.
\end{prop}

\begin{proof}
Let $a_{0}=0$, and let $(a_{n})_{n\in I}$ be the positive atoms of $\lambda$ arranged in some arbitrary order; here, $I=\n$ if there is an infinite number of atoms and $I=[N]$ for some $N$ otherwise. For every $n\in I$, denote $p_{n}=\theta(\{ a_{n}\})$, and define $q_{n}$ by
\[
q_{n}=\frac{p_{n}}{p_{0}+p_{1}+\ldots+p_{n}}.
\]
For $n\in I$, let $(S_{n},d_{n},\mu_{n})$ be the metric space corresponding to the Bernoulli measure $\theta_{n}=(1-q_{n})\delta_{0}+q_{n}\delta_{a_{n}}$, obtained by multiplying the metric described in the examples section by $a_{n}$. Then the assumption in the construction of the selection space and the two conditions in Proposition~\ref{prop:selection_space_properties} hold for $\{ (S_{n},d_{n},\mu_{n})\} _{n\in I}$:
\begin{enumerate}
\item When constructing a Bernoulli random variable $p\delta_{0}+(1-p)\delta_{1}$with probability $p$ of obtaining $0$ as in Example~\ref{enu:bernoulli_construction}, the parameters $m$ and $\alpha$ may be chosen so that $\alpha\geq p$. Indeed, setting $f_{m}(x)=x^{2}+(\frac{1-x}{m-1})^{2}$, we have $p=f_{m}(\alpha)$. Since $\lim_{m\to\infty}f_{m}(p)=p^{2}<p$, there exists an $m$ large enough so that $f_{m}(p)<p$, and by continuity of $f_{m}$ and the fact that $f_{m}(1)=1$, there exists $\alpha\in[p,1]$ with $f_{m}(\alpha)=p$ as required. This shows that the probability to obtain $0$ under $\mu_{n}$ is bounded by
\begin{equation}\label{eq:bound_on_q}
\mu_{n}\left(\left\{ 0\right\} \right)\geq1-q_{n}.
\end{equation}
Thus
\begin{equation}\label{eq:sum_of_q_is_finite}
\sum_{n\in I}1-\mu_{n}\left(\left\{ 0\right\} \right)\leq\sum_{n\in I}q_{n}=\sum_{n\in I}\frac{p_{n}}{p_{0}+p_{1}+\ldots+p_{n}}\leq\sum_{n\in I}\frac{p_{n}}{p_{0}}=\frac{1-p_{0}}{p_{0}},
\end{equation}
and the sum in~\eqref{eq:prob_sums_are_finite} is finite.
\item Since $a_{n}\in[a,2a]$, the requirement~\eqref{eq:metrics_are_subadditive_assumption} is immediate: the left-hand side of~\eqref{eq:metrics_are_subadditive_assumption} is at most $2a$, while the right-hand side is at least $2a$.
\item $\inf I=1>-\infty$ as required in~\eqref{eq:assumption_separability}.
\end{enumerate}
Thus, the selection space $(S,d,\mu)$ constructed from $\{(S_{n},d_{n},\mu_{n})\} _{n\in I}$ is a complete separable measured metric space. Since the positive atoms of the $\mu_{n}$s are all distinct, equations~\eqref{eq:prob_of_0} and~\eqref{eq:prob_of_positive} imply that
\begin{equation}\label{eq:single_prob_zero}
\p\left[d\left(X,Y\right)=0\right]=\prod_{n\in I}\p\left[d_{n}\left(X_{n},Y_{n}\right)=0\right]=\prod_{n\in I}\left(1-q_{n}\right)=p_{0},
\end{equation}
and
\[
\p\left[d\left(X,Y\right)=a_{n}\right]=q_{n}\prod_{i\in I,\,i\,>\,n}\left(1-q_{i}\right)=p_{n}
\]
as needed.

We remark that by~\eqref{eq:bound_on_q} and~\eqref{eq:single_prob_zero},
\begin{equation}\label{eq:entire_vector_is_0_whp}
\p\left[X=\left(0,0,\ldots\right)\right]=\prod_{n\in I}\mu_{n}\left(\left\{ 0\right\} \right)\geq\prod_{n\in I}\left(1-q_{n}\right)=p_{0}.\qedhere
\end{equation}
\end{proof}
\begin{proof}[Proof of Theorem~\ref{thm:discrete_distributions}]\label{prfTheo1.1}
The measure $\theta$ can be written as
\[
\theta=p_{\infty}\delta_{0}+\sum_{n\in I}p_{n}\lambda_{n},
\]
where $I\subseteq\z$, and for every $n\in I$, $p_{n}>0$ and the measure $\lambda_{n}$ is discrete and supported on $(2^{n},2^{n+1}]$. Note that if $p_{\infty}=0$, then necessarily $\inf I=-\infty$. Define the numbers
\[
\beta_{n}=\frac{p_{n}}{1-\sum_{i\in I,\,i\,>\,n}p_{i}},
\]
and let $(S_{n},d_{n},\mu_{n})$ be the metric space defined in the proof of Proposition~\ref{prop:positive_interval_is_feasible} corresponding to the measure $\theta_{n}=(1-\beta_{n})\delta_{0}+\beta_{n}\lambda_{n}$. Note that each $S_{n}$ is countable. Then the assumption in the construction of the selection space and the two conditions in Proposition~\ref{prop:selection_space_properties} hold for $\{(S_{n},d_{n},\mu_{n})\} _{n\in I}$:
\begin{enumerate}
\item If $\sup I<\infty$ then of course the sum in~\eqref{eq:prob_sums_are_finite} is finite. Otherwise, by~\eqref{eq:entire_vector_is_0_whp}, if $X_{n}\sim\mu_{n}$, then
\[
\p\left[X_{n}=\left(0,0,\ldots\right)\right]\geq1-\beta_{n}.
\]
Thus
\begin{align*}
\sum_{n\in I\,\intersect\,\n} \p\left[X_{n}\neq\left(0,0,\ldots\right)\right]&\leq\sum_{n\in I\,\intersect\,\n}1-\left(1-\beta_{n}\right) =\sum_{n\in I\,\intersect\,\n}\beta_{n} \\ &\leq\sum_{n\in I\,\intersect\,\n}\frac{p_{n}}{1-\sum_{i\in I,\,i\,>\,n}p_{i}}\leq\frac{1}{p_{\infty}+\sum_{i\in I,\,i\,<\,1}p_{i}}.
\end{align*}
\item Since the range of $d_{n}$ is $(2^{n},2^{n+1}]$, the requirement~\eqref{eq:metrics_are_subadditive_assumption} is immediate; in fact, the left hand side of~\eqref{eq:metrics_are_subadditive_assumption} is bounded by just half of the right hand side.
\item $\lim_{n\to-\infty}\sup_{z,w\in S_{n}}d_{n}(z,w)=\lim_{n\to-\infty}2^{n+1}=0$ as required in~\eqref{eq:assumption_separability}.
\end{enumerate}
Thus, the selection space $(S,d,\mu)$ constructed from $\{(S_{n},d_{n},\mu_{n})\} _{n\in I}$ is a complete separable measured metric space. Since the supports of $\lambda_{n}$ are disjoint, equations~\eqref{eq:prob_of_0} and~\eqref{eq:prob_of_positive} imply that
\[
\p\left[d\left(X,Y\right)=0\right]=\prod_{n\in I}\p\left[X_{n}=Y_{n}\right]=\prod_{n\in I}\left(1-\beta_{n}\right)=p_{\infty},
\]
while the probability to sample from $\lambda_{n}$ is exactly $\beta_{n}\prod_{i\in I,\,i\,>\,n}(1-\beta_{i})=p_{n}$, as needed.
\end{proof}

\section{Proof of Theorem~\ref{thm:bounded_continuous_distributions}}

\hypertarget{target:proof_of_continuous}Our proof first finds an easy-to-construct feasible ``starter'' distribution, and then changes its metric in order to obtain the distribution $\theta$. The idea is as follows. Let $W$ be a random variable on $[0,1]$ whose cumulative distribution function $F_{W}$ is continuous and strictly increasing on $[0,1]$, and let $\varphi:[0,1]\to[0,1]$ be a strictly increasing bijection. The cumulative distribution function of the random variable $\varphi(W)$ is given by
\[
F_{\varphi\left(W\right)}\left(t\right)=\p\left[\varphi\left(W\right)\leq t\right]=\p\left[W\leq\varphi^{-1}\left(t\right)\right]=F_{W}\left(\varphi^{-1}\left(t\right)\right).
\]
Denoting by $G$ the cumulative distribution function of $\theta$ and choosing $\varphi(t)=G^{-1}(F_{W}(t))$ gives us, for $t\in[0,1]$, $F_{W}(\varphi^{-1}(t))=G(t)$, and so the random variable $\varphi(W)$ has distribution $\theta$. If $W$ is achievable by some $(S,d,\mu)$, then when $X,Y\sim\mu$ are independent, we have $\varphi(d(X,Y))\sim\theta$. However, this does not immediately mean that $(S,\varphi\circ d,\mu)$ is a measured metric space that achieves $\theta$, since there is no guarantee that $\varphi\circ d$ is still a metric (or that $\mu$ is still Borel under the topology induced by $\varphi\circ d$). The essence of the proof is to find a suitable $W$ which allows $\varphi$ to preserve the metric $d$.
\begin{defi}
A function $f:\nonnegative\to\nonnegative$ is called \emph{metric preserving }if for every metric space $(S,d)$, the function $f\circ d$ is also a metric on $S$. It is called \emph{strongly metric preserving} if it is metric preserving and $(S,d)$ is topologically equivalent to $(S,f\circ d)$.
\end{defi}

Metric preserving functions have been well studied; see~\cite{corazza_introduction_to_metric_preserving_functions} for a brief introduction. We will require only the following standard result.

\begin{theo}[{\cite[p.~131]{kelley_general_topology}}]\label{thm:subadditivity_implies_metric_preserving}
Let $f:\nonnegative\to\nonnegative$ be continuous, nondecreasing, and satisfying the following two conditions:
\begin{enumerate}
\item $f(x)=0\iff x=0$.
\item (Subadditivity) $f(x+y)\leq f(x)+f(y)$ for all $x,y\in\nonnegative$.
\end{enumerate}
Then $f$ is strongly metric preserving.
\end{theo}

Theorem~\ref{thm:subadditivity_implies_metric_preserving} implies that it suffices to find a continuous random variable $W$ such that $G^{-1}(F_{W}(t))$ is subadditive. If $G^{-1}$ is itself already subadditive, then no modification by $F_{W}$ is actually needed; in this case, we can take $W$ to be the uniform measure on $[0,1]$, so that $F_{W}(t)=t$ for $t\in[0,1]$ and $G^{-1}(F_{W}(t))=G^{-1}(t)$. Obtaining $\theta$ is then just an instance of the inverse sampling theorem (as we saw in the examples section, the uniform measure is indeed feasible). However, when $G^{-1}$ is not subadditive, the idea is that if $F_{W}$ itself is in some sense very strongly subadditive, then its subadditivity is enough to counter the non-subadditivity of~$G^{-1}$. This is made precise in the following simple proposition.
\begin{prop}\label{prop:composing_into_subadditivity}
Let $\Psi:[0,1]\to[0,1]$ be an absolutely continuous nondecreasing function, and assume that there exist $m,M>0$ such that
\begin{equation}\label{eq:h_has_pinched_derivatives}
m<\Psi'\left(t\right)<M.
\end{equation}
Let $f:\nonnegative\to[0,1]$ be such that for every $x,y\in\nonnegative$ with $x\leq y$,
\begin{equation}\label{eq:f_is_strongly_subadditive}
f\left(x+y\right)\leq\frac{m}{M}f\left(x\right)+f\left(y\right).
\end{equation}
Then $\Psi\circ f$ is subadditive.
\end{prop}

\begin{proof}
By absolute continuity, for all $x\leq y$,
\begin{align*}
\Psi\left(f\left(x+y\right)\right) & =\Psi\left(f\left(y\right)\right)+\int_{f\left(y\right)}^{f\left(x+y\right)}\Psi'\left(t\right)\,d t\\
\overset{\eqref{eq:h_has_pinched_derivatives}}&\leq\Psi\left(f\left(y\right)\right)+M\left(f\left(x+y\right)-f\left(y\right)\right)\\
\overset{\eqref{eq:f_is_strongly_subadditive}}&\leq\Psi\left(f\left(y\right)\right)+mf\left(x\right)
\\
 \overset{\eqref{eq:h_has_pinched_derivatives}}&\leq\Psi\left(f\left(y\right)\right)+\Psi\left(f\left(x\right)\right).\qedhere
\end{align*}
\end{proof}

The next proposition shows that it is indeed possible to obtain such $f$.
\begin{prop}\label{prop:strongly_subadditive_functions_are_realizable}
For every $\eps\in(0,1)$, there exists a feasible continuous random variable $W$ whose cumulative distribution function $F_{W}$ satisfies
\[
F_{W}\left(x+y\right)\le\eps F_{W}\left(x\right)+F_{W}\left(y\right)
\]
for all $x\leq y$.
\end{prop}

Given the above, Theorem~\ref{thm:bounded_continuous_distributions} quickly follows.
\begin{proof}[Proof of Theorem~\ref{thm:bounded_continuous_distributions}]
By~\eqref{eq:bounds_on_density_function}, $G^{-1}$ is absolutely continuous, strictly increasing, and has upper and lower bounds on its derivative $(G^{-1})(t)'=\frac{1}{g(G^{-1}(t))}$. Then by Propositions~\ref{prop:composing_into_subadditivity} and~\ref{prop:strongly_subadditive_functions_are_realizable}, there is a feasible $W$ achieved by $(S,d,\mu)$ such that $\varphi=G^{-1}\circ F_{W}$ is subadditive. Since both $G^{-1}$ and $F_{W}$ are continuous, by Theorem~\ref{thm:subadditivity_implies_metric_preserving}, $\varphi$ is strongly metric preserving, and so $(S,\varphi\circ d,\mu)$ achieves $\theta$.
\end{proof}
\begin{proof}[Proof of Proposition~\ref{prop:strongly_subadditive_functions_are_realizable}]
Let $\alpha\in(0,\frac{1}{2})$ to be chosen later, define $H:\r\to[0,1]$ by
\[
H\left(t\right)=
\begin{cases}
0 & t<0\\
t^{\alpha} & 0\leq t\leq1\\
1 & t>1,
\end{cases}
\]
and let $\mu$ be the distribution whose cumulative distribution function is $H$. This distribution has density $h(t)=\alpha t^{\alpha-1}$ for $t\in[0,1]$. Let $W=\abs{X-Y}$, where $X,Y\sim\mu$ are i.i.d. By definition, $W$ is feasible. For $t\in[0,1]$, the cumulative distribution function $F_{W}$ is given by
\begin{align*}
F_{W}\left(t\right) & =\int_{-\infty}^{\infty}\int_{-\infty}^{\infty}h\left(x\right)h\left(y\right)\one_{\abs{x-y}\,\leq\,t}\,dydx\\
& =2\int_{0}^{1}\int_{x}^{x+t}h\left(x\right)h\left(y\right)\,dydx\\
& =2\int_{0}^{1}h\left(x\right)\left(H\left(x+t\right)-H\left(x\right)\right)\,dx\\
& =2\int_{0}^{1-t}\alpha x^{2\alpha-1}\left(\left(1+\frac{t}{x}\right)^{\alpha}-1\right)\,dx+2\int_{1-t}^{1}\alpha x^{\alpha-1}\left(1-x^{\alpha}\right)\,dx.
\end{align*}
The function $F_{W}$ can readily be shown to be concave on $\nonnegative$ by calculating the second derivative of $2\int_{0}^{1}\int_{x}^{x+t}h(x)h(y)dydx$. Thus, for all $z,w\in\nonnegative$,
\[
F_{W}\left(z\right)\leq F_{W}\left(w\right)+\left(z-w\right)F_{W}'\left(w\right).
\]
Let $x,y\in(0,\infty)$ with $x\leq y$. Choosing $z=x+y$ and $w=y$, we get
\begin{align*}
F_{W}\left(x+y\right) & \leq F_{W}\left(y\right)+xF_{W}'\left(y\right)\\
& \leq F_{W}\left(y\right)+xF_{W}'\left(x\right)\\
& =F_{W}\left(y\right)+F_{W}\left(x\right)\cdot\frac{xF_{W}'\left(x\right)}{F_{W}\left(x\right)}.
\end{align*}
It therefore suffices to show that for all $t\in\nonnegative$,
\[
\frac{tF_{W}'\left(t\right)}{F_{W}\left(t\right)}\leq\eps.
\]
As $F_{W}$ is concave and increasing, it is enough to show a bound of $\frac{1}{2}\eps$ for all $t\leq\frac{1}{2}$, since for all $t>\frac{1}{2}$, $F_{W}'(t)\leq F_{W}'(\frac{1}{2})$ and $F_{W}(t)\geq F_{W}(\frac{1}{2})$. Assuming that $t\leq\frac{1}{2}$, the denominator is bounded by
\begin{align*}
F_{W}\left(t\right) & =2\alpha\left(\int_{0}^{1-t}x^{2\alpha-1}\left(\left(1+\frac{t}{x}\right)^{\alpha}-1\right)\,dx+\int_{1-t}^{1}x^{\alpha-1}\left(1-x^{\alpha}\right)\,dx\right)\\
& \geq2\alpha\int_{0}^{t}x^{2\alpha-1}\left(\left(\frac{t}{x}\right)^{\alpha}-1\right)\,dx=2\alpha\left(\frac{1}{\alpha}t^{2\alpha}-\frac{1}{2\alpha}t^{2\alpha}\right)=t^{2\alpha}.
\end{align*}
Recalling that $\alpha<\frac{1}{2}$, the numerator is bounded by
\begin{align*}
F_{W}'\left(t\right) & =2\int_{0}^{1}h\left(x\right)h\left(x+t\right)\,dx\\
& =2\alpha^{2}\left(\int_{0}^{t}x^{2\alpha-2}\left(1+\frac{t}{x}\right)^{\alpha-1}\,dx+\int_{t}^{1-t}x^{2\alpha-2}\left(1+\frac{t}{x}\right)^{\alpha-1}\,dx\right)\\
& \leq2\alpha^{2}\left(\int_{0}^{t}x^{2\alpha-2}\left(\frac{t}{x}\right)^{\alpha-1}\,dx+\int_{t}^{1-t}x^{2\alpha-2}\,dx\right)\\
& =2\alpha t^{2\alpha-1}-2\alpha^{2}\frac{1}{1-2\alpha}\left(\frac{1}{\left(1-t\right)^{1-2\alpha}}-\frac{1}{t^{1-2\alpha}}\right)\\
& \leq2\alpha\frac{1}{t^{1-2\alpha}}\left(1+\frac{\alpha}{1-2\alpha}\right).
\end{align*}
Choosing $\alpha\leq\frac{1}{4}$ then gives $F_{W}'(t)\leq4\alpha t^{2\alpha-1}$. We thus have
\[
\frac{tF_{W}'\left(t\right)}{F_{W}\left(t\right)}\leq\frac{4\alpha t^{2\alpha}}{t^{2\alpha}}=4\alpha,
\]
and the proposition follows for $\alpha\leq\frac{1}{8}\eps$.
\end{proof}
\begin{rema}
The function $F_{W}$ may be expressed in a slightly more closed form using hypergeometric functions; a short calculation reveals that for $t\in[0,1]$,
\[
F_{W}\left(t\right)=1+2\left(1-t\right)^{\alpha}\left(\hypergeometric{-\alpha}1{1+\alpha}{1-t}-1\right),
\]
where $\hypergeometric abcz$ is the hypergeometric function
\[
\hypergeometric abcz=\frac{\Gamma\left(c\right)}{\Gamma\left(b\right)\Gamma\left(c-b\right)}\int_{0}^{1}\frac{x^{b-1}\left(1-x\right)^{c-b-1}}{\left(1-xz\right)^{a}}\,dx.
\]
A plot of $F_{W}$ for some values of $\alpha$ is given in Figure~\ref{fig:F_W}, where the asymptotic behavior of $t^{2\alpha}$ can be seen near the origin.

\begin{figure}[h]
\begin{centering}
\includegraphics[scale=0.7]{feasible_fw}
\end{centering}
\caption{The function $F_{W}$ for various values of $\alpha$.\label{fig:F_W}
}
\end{figure}
\end{rema}

\begin{rema}
It might perhaps be possible to extend the proof method of Theorem~\ref{thm:bounded_continuous_distributions} to distributions on $\nonnegative$ where the density is always positive. Indeed, if we compose $G^{-1}$ with another function $f$, we still have
\begin{align*}
G^{-1}\left(f\left(x+y\right)\right) & =G^{-1}\left(f\left(y\right)\right)+\int_{f\left(y\right)}^{f\left(x+y\right)}\frac{d}{d t}G^{-1}\left(t\right)\,d t\\
& \leq G^{-1}\left(f\left(y\right)\right)+\left(f\left(x+y\right)-f\left(y\right)\right)\sup_{t\in\left[f\left(y\right),\,f\left(x+y\right)\right]}\frac{d}{d t}G^{-1}\left(t\right).
\end{align*}
An analogue of Theorem~\ref{thm:bounded_continuous_distributions} would follow if we could find, for each $\eps>0$, an unbounded feasible random variable $W$ such that for all $x\leq y$
\[
F_{W}\left(x+y\right)\leq F_{W}\left(y\right)+\left(\eps\sup_{t\in\left[F_{W}\left(y\right),\,F_{W}\left(x+y\right)\right]}\frac{d}{d t}G^{-1}\left(t\right)\right)F_{W}\left(x\right).
\]
Unlike the random variable from Proposition~\ref{prop:strongly_subadditive_functions_are_realizable}, such a $W$ would have to rely directly on $G$ in some way. However, the method cannot be easily extended to distributions whose density is $0$ in an interval: this would imply that $G^{-1}\circ F_{W}$ is discontinuous, which would prevent it from preserving the metric.
\end{rema}


\bibliography{gross}
\end{document}
